\begin{figure}[ht]
\centering
\begin{tikzpicture}[x=1cm, y=1cm, font=\small]

% Axes: harmonic amplitude (relative to fundamental) versus harmonic number
\draw[->, thick] (-0.2, 0) -- (9.2, 0) node[right, font=\footnotesize] {harmonic $n$};
\draw[->, thick] (0, -0.1) -- (0, 5.4) node[above, font=\footnotesize] {relative amplitude};

% y grid: amplitude scaled so 1.0 = 5 cm
\foreach \a/\lab in {0/0, 2.5/0.5, 5.0/1.0} {
  \draw[thick] (-0.1, \a) -- (0.1, \a);
  \node[anchor=east, font=\footnotesize] at (-0.1, \a) {\lab};
}

% Harmonics 1..7 at x positions; three waveform families offset within each group
% square wave: odd harmonics only, amplitude 1/n
% sawtooth: all harmonics, amplitude 1/n
% triangle: odd harmonics only, amplitude 1/n^2
\foreach \n in {1,2,3,4,5,6,7} {
  \pgfmathsetmacro{\xc}{1.1*\n}
  \node[anchor=north, font=\footnotesize] at (\xc, -0.05) {\n};
  % sawtooth (all n): 1/n, amber
  \pgfmathsetmacro{\saw}{5.0/\n}
  \fill[engamber] ({\xc-0.30}, 0) rectangle ({\xc-0.12}, \saw);
  % square (odd n): 1/n, blue
  \pgfmathsetmacro{\isodd}{mod(\n,2)}
  \ifdim\isodd pt>0.5pt
    \pgfmathsetmacro{\sq}{5.0/\n}
    \fill[engblue] ({\xc-0.10}, 0) rectangle ({\xc+0.08}, \sq);
    % triangle (odd n): 1/n^2, red
    \pgfmathsetmacro{\tri}{5.0/(\n*\n)}
    \fill[engred] ({\xc+0.10}, 0) rectangle ({\xc+0.28}, \tri);
  \fi
}

% Legend
\fill[engamber] (5.6, 4.9) rectangle (5.9, 5.1);
\node[anchor=west, font=\footnotesize] at (6.0, 5.0) {sawtooth ($1/n$, all $n$)};
\fill[engblue] (5.6, 4.4) rectangle (5.9, 4.6);
\node[anchor=west, font=\footnotesize] at (6.0, 4.5) {square ($1/n$, odd $n$)};
\fill[engred] (5.6, 3.9) rectangle (5.9, 4.1);
\node[anchor=west, font=\footnotesize] at (6.0, 4.0) {triangle ($1/n^2$, odd $n$)};

\end{tikzpicture}
\caption[Harmonic amplitudes of three canonical waveforms]{Relative
harmonic amplitudes for the sawtooth ($1/n$, every integer harmonic),
the square wave ($1/n$, odd harmonics only), and the triangle wave
($1/n^2$, odd harmonics only), each normalised to a unit fundamental.
The triangle's $1/n^2$ decay is visibly faster than the square's
$1/n$; the smoother time-domain waveform carries the more compact
spectrum. The even harmonics vanish for the square and triangle
because both are odd-symmetric about their half-period.}
\label{fig:vol01:ch06:harmonic-content}
\end{figure}
